% Everything the document needs before the first section: packages, page setup,
% the theorem and pseudocode environments, the drawing styles shared by every
% figure, and the notation macros. Sections and figures declare nothing of their
% own, so one edit here restyles the whole report.

\usepackage[T1]{fontenc}
\usepackage[utf8]{inputenc}
\usepackage{lmodern}
\usepackage[english]{babel}
\usepackage{microtype}
\usepackage{geometry}
\geometry{margin=2.4cm}

\usepackage{amsmath,amssymb,amsthm}
\usepackage{booktabs}
\usepackage{enumitem}
\usepackage{caption}
\usepackage{subcaption}
\usepackage{xcolor}
\usepackage[linesnumbered,ruled,vlined]{algorithm2e}
\usepackage{tikz}
\usetikzlibrary{arrows.meta,calc,fit,positioning,backgrounds}
\usepackage{import}
\usepackage{hyperref}
\usepackage{bookmark}

% --- Shared palette ------------------------------------------------------
% Named by the role each colour plays in a figure, not by the colour itself,
% so a change of palette does not turn every caption into a lie.
\definecolor{linkblue}{HTML}{1F4E79}
\definecolor{doorgreen}{HTML}{2E7D32}   % an edge kept: a door
\definecolor{wallred}{HTML}{C62828}     % an edge deleted: a wall segment
\definecolor{cellgrey}{HTML}{ECEFF1}    % the area a region covers
\definecolor{exitorange}{HTML}{E65100}  % the single exit cell

\hypersetup{colorlinks=true,linkcolor=linkblue,citecolor=linkblue,urlcolor=linkblue}

% --- Shared drawing styles ----------------------------------------------
\tikzset{
  cell/.style     = {draw=black!45, fill=cellgrey, minimum size=9mm},
  node vertex/.style = {circle, draw=black!70, fill=white, inner sep=0pt,
                        minimum size=4.2mm, font=\scriptsize},
  exit vertex/.style = {node vertex, draw=exitorange, very thick,
                        fill=exitorange!12},
  door/.style     = {doorgreen, line width=1.1pt},
  wall/.style     = {wallred, line width=1.1pt, dash pattern=on 2pt off 2pt},
  wall segment/.style = {wallred, line width=1.6pt},
  region/.style   = {draw=black!60, fill=cellgrey, thick},
  link/.style     = {-{Stealth[length=2mm]}, black!65},
  tree node/.style = {draw=black!70, rounded corners=1.5pt, inner sep=2.5pt,
                      font=\scriptsize, fill=white},
  tree leaf/.style = {tree node, draw=doorgreen!80!black, fill=doorgreen!8},
  tree edge/.style = {black!55},
}

% --- Theorem-like environments ------------------------------------------
\theoremstyle{definition}
\newtheorem{definition}{Definition}
\theoremstyle{plain}
\newtheorem{proposition}{Proposition}
\theoremstyle{remark}
\newtheorem*{remark}{Remark}

% --- Notation ------------------------------------------------------------
% The report says "the grid graph" and "the maze" constantly; these keep the
% symbols identical in prose, in displays and in pseudocode.
\newcommand{\gridgraph}{G_{\text{grid}}}
\newcommand{\gridedges}{E_{\text{grid}}}
\newcommand{\maze}{M}
\newcommand{\cells}{V}
\newcommand{\region}[2]{[#1]\times[#2]}
\newcommand{\exitcell}{v_{\text{exit}}}
\SetKwProg{Fn}{Function}{:}{}
